\subsection{The actual selected-prime theta chain}
\label{sec:actual-theta-chain}

The ordinary CM construction gives normalized theta functions and their
simultaneous finite moment laws.  We now identify the resulting local
operators with the original arithmetic chain, including its point
models, integral coefficient lattices, and twisted initial source.
The phases in this identification are retained, rather than replaced by
their absolute values.  The source interfaces are attributed to
Alp\"oge's pinned construction.\cite{theta-chain-source}

Fix one of its admissible finite column data, a word prime $T$, a column
index, and the common coefficient field $K$ already chosen for its
characters, normalized theta jets and direction constants.  Write
$E$ for the CM field, $[E:\mathbb Q]=2f$, and
\[
 h=(t+1)/2,\qquad N_i=k+i(t+1)\quad(0\le i\le m).
\]
Here $t$ is odd and the $N_i$ are even, as in the original datum.
The selected primes over $T$ have residue field $\mathbb F_T$, are
distinct from their conjugates and from the auxiliary levels and the
different.  Retain the original coefficient weights, all character
values, the primitive theta normalization and the differential
directions.  The numerical admissibility includes the original hook
bound at each target and $T>2fN_m+1$.  It also includes the original
coefficient weight inequalities used below.  These are conditions on
the admitted column datum, not new estimates on a varying field.

Let $\mathscr S_i(w)$ be the original local integral lattice of
equivariant class functions, embedded in its original complex space by
$\sigma:K\hookrightarrow\mathbb C$.  Write $K_i$ for the original
class operator from stage $i-1$ to stage $i$, with its original
automorphism masses, and $C_m=K_m\cdots K_1$.  At $w\mid T$ of
selected type, let $J$ be the prescribed selected embedding set and
let $U_i$ be the ordered product of its local neighbor operators.

\begin{theorem}[Actual local realization and initial-source compatibility]
\label{thm:actual-theta-chain}
For the finite column datum just specified, the normalized CM theta
construction supplies the original point kernels and their original
integral value-lattice identifications.  The kernels agree with the
literal class maps $K_i$, with scalar one.  They satisfy the original
target equivariance, simultaneous dilation and Hecke laws, and both
selected and conjugate integral kernel laws.

In particular, put
\[
 d_i=fN_{i-1}+\sum_\tau s_{i-1,\tau}
               +\sum_{\tau\notin J}r_{i,\tau}.
\]
Here $s_{i-1,\tau}$ is the original determinant shift and
$r_{i,\tau}$ is the first part of the original residual target
partition in the allowance-removal calculation.  Then
\begin{align}
 T^{d_i}U_iK_i\mathscr S_{i-1}(w)&\subseteq\mathscr S_i(w),
       &U_iK_i&=T^{h|J|}K_iU_{i-1}.
       \label{eq:actual-chain-step}
\end{align}
The original twisted initial lattice $\mathscr N(w)$ is the
\emph{rational image} of
\[
 T^aU_0^m\mathscr S_{\rm ram}(w),\qquad
 a=m\sum_{\tau\in J}b_\tau,\qquad
 \mathscr S_{\rm ram}(w)\subseteq\mathscr S_0(w),
\]
under the original initial class identification.  Consequently
\[
 T^\ell C_m\mathscr N(w)\subseteq\mathscr S_m(w),\qquad
 \ell=\sum_{i=1}^m d_i+h|J|\frac{m(m+1)}2-a.
 \label{eq:actual-chain-selected}
\]
At a conjugate place the same actual maps satisfy
\[
 K_i\mathscr S_{i-1}(w)\subseteq T^{fh}\mathscr S_i(w),
 \qquad C_m\mathscr S_0(w)\subseteq T^{mfh}\mathscr S_m(w).
 \label{eq:actual-chain-conjugate}
\]
All inclusions are assertions about the prescribed embedded lattices;
they construct their unique rational target lifts.  The point-model
Hermitian form is exactly the original frame form, and the class form
has the original automorphism masses.  No normalization of these forms
or change of coefficient field is made.
\end{theorem}

\begin{proof}
We give the construction in five steps.

\emph{1. Points, phases and local modifications.}
At rank $N_i$, let $d$ denote one of the original decorated lattices:
its self-dual Hermitian lattice, selected residue functionals,
prime-to-$T$ level structures, and the additional fixed markings.
At a selected prime $P$ record the pair
$(N_P,\ell_P)$, where $\ell_P:N_P\to\mathbb F_T$ is the original
nonzero functional.  Record all the remaining data as an off-$T$
class.  Retain a phase $\zeta\in\mu(K)$ as well.

More precisely, within an original component $b$, identify
$(\zeta,d)$ and $(\zeta',d')$ in the off-$T$ record if an original
change of markings $g$ has underlying rational unitary map $1$,
if $\zeta=\zeta'\chi_b(g)$, and if $gd,d'$ agree off $T$.
Let $Z_i$ be the disjoint union of these quotients and set
\[
 P_i=Z_i\times\prod_{P\mid T\ {
m selected}}
           \{(N_P,\ell_P)\}.
\]
Every such point has a representative $(b,\zeta,d)$.  To see this,
prescribe the displayed lattices at the selected primes and their
Hermitian duals at the conjugate primes, keeping the old lattices
elsewhere.  The intersection of these finitely many local conditions
is a global lattice with exactly those localizations; duality gives
self-duality.  The residue functionals lift by the Chinese remainder
theorem.  All other markings transport through unchanged localizations,
because their levels are prime to $T$.  The same construction makes
one prescribed upper or lower neighbor, or one residue-functional
dilation, leaving every other local component unchanged.

Equality of two representatives as full points is therefore exactly
\[
 \exists g:\quad \gamma(g)=1,\quad
 \zeta=\zeta'\chi_b(g),\quad gd=d'.                 \tag{1}
\]
Indeed the off-$T$ equality and equality of the selected pairs imply
equality of the global decorated lattices by localization; the reverse
implication is immediate.  This retains the original level markings,
not just the underlying lattices.

The group $G_i=U(E^{N_i})\times\mu(K)$ acts as follows.  If
$g$ is any original lift of $\gamma\in U(E^{N_i})$, define
\[
 (\gamma,\eta)(b,\zeta,d)
   =(b,\eta\zeta\chi_b(g)^{-1},gd).                 \tag{2}
\]
The original marking group surjects onto the rational unitary group:
act by that unitary map on the lattice and transport the markings.
Equation (1) makes (2) independent of the lift and representative;
the character multiplication law proves the group law.  Its orbits
are exactly the original classes.  On a stabilizer the map
$g\mapsto(\gamma(g),\chi_b(g))$ is bijective: surjectivity follows
from (1), and injectivity follows because an automorphism fixing the
decorated lattice with underlying map $1$ fixes every marking.
Thus the stabilizer orders, and hence all class masses, are unchanged.

\emph{2. Coefficients and their integral lattices.}
Let $V_i(K)$ be the original highest-weight polynomial space with its
determinant twist.  The actual point value space is its evaluation on
the original frames, extended to $\mathbb C$.  Evaluation is injective:
after multiplying by the fixed determinant character, a polynomial
vanishing on all invertible frames vanishes on a Zariski-dense open
set.  Its image spans the value space by its definition.  Consequently
evaluation has an inverse coordinate map on that space.  In these
coordinates the stage-$i$ weight is the original frame weight shifted
by $ih$ in every entry.  This is exactly the source's shift
$i(t+1)/2$, since $t+1=2h$; the underlying polynomial partition is
unchanged.  This observation retains the determinant twist in the
coordinate identification.

At $w\mid T$, use the embedding $\tau:E\to K$ in selected type
and $\tau\circ\overline{\phantom{x}}$ in conjugate type.
For a subset $A$ of selected primes, the allowed local lattice is
$N_P$ in selected type at $P\notin A$, and $\ker\ell_P$ otherwise.
Choose a local basis of each such lattice and act on the original
integral highest-weight lattice by its exact determinant-twisted
polynomial action.  Changing the basis by $\operatorname{GL}_{N_i}$
over the local ring does not change this lattice: the action and its
inverse are integral.  In conjugate type the frame is the dual frame,
so this construction is the original coefficient lattice on the
Hermitian dual of the patched lattice.  In selected type it is the
original coefficient lattice on the patched lattice itself.  In
particular $A=\varnothing$ gives the original source value lattice,
$A$ equal to all selected primes gives its fully allowed kernel
target, and $A=\{P\}$ gives the original single-prime allowance.
There is no intervening index or saturation.

Off $T$, define the value lattice at an off-$T$ point by requiring
membership in the original transported coefficient lattice for
\emph{every} presentation of that point.  This is precisely the
intersection of those transported lattices.  Its invariance follows
by reindexing the presentations.  The root-of-unity factor in (2) is
an integral unit at every finite place.

These pointwise conditions characterize the embedded original local
class lattice, in both directions.  For the nontrivial direction,
recover from each integral point value its unique $K$-polynomial
using the coordinate map just constructed.  Equivariance over
$\mathbb C$ and injectivity of evaluation imply the original
stabilizer equations over $K$.  There are finitely many classes.
For each class choose a coefficient denominator outside $w$; their
product remains outside $w$ and clears all values simultaneously.
The resulting rational equivariant class function belongs to the
original localized lattice.  This proves rationality of the lift;
it does not assert that an arbitrary complex vector is rational.

The Hermitian form used throughout is the original frame form
$\operatorname{ipVal}$, restricted to its value space.
The action (2) on values is the original frame action multiplied by
$\sigma(\eta)$.  Unitarity of $\gamma$ preserves that frame form
and $|\sigma(\eta)|=1$.  Thus this construction keeps its exact form;
it does not compare it with a separately chosen coefficient norm.

\emph{3. The same theta kernel, including its phase.}
For an actual pair $d,d'$, form the finite indicator on the original
Hom-lattice cosets.  Apply the original normalized theta family and
its prescribed mixed derivatives in the directions $c_{D,\tau}$.
With source weight $\lambda$ and the original nonnegative determinant
exponents $n_\tau$, expand
\[
 F(X)\prod_\tau\det(X_\tau)^{n_\tau}
       \bigg|_{X_\tau=Z_\tau Y_{\tau,[1,N_{i-1}]}}
       =\sum_\alpha Z^\alpha Q_\alpha(F)(Y).
\]
The actual coefficient kernel is
\[
 k_i(d,d')F=\eta_i(d,d')\sum_\alpha J_\alpha(d,d')Q_\alpha(F).
                                                               \tag{3}
\]
Here $J_\alpha$ is the finite indicator sum of the original ordinary
mixed derivatives; it is not a divided Taylor coefficient.
The previously established normalized CM construction supplies these
same $J_\alpha$ and their simultaneous moment laws.  Here is the
target-representation check, including its determinant shifts.
At one embedding put $a=N_{i-1}<b=N_i$ and
$\nu_j=\lambda_j-\lambda_{a-1}$.  The original weight recursion is
\[
 n=\lambda_{a-1}+s^{\sharp},\qquad
 \lambda'_j=\begin{cases}
       \lambda_j+s^{\sharp}-s^{\rm div}&j<a,\\
       -s^{\rm div}&j\ge a.
                 \end{cases}
\]
Consequently the target residual partition is
$(\nu_0+n,\ldots,\nu_{a-1}+n,0,\ldots,0)$.
For any upper triangular $B\in\operatorname{GL}_b$,
$(YB)_{[1,a]}=Y_{[1,a]}B_{[1,a],[1,a]}$.  The original
column-highest identity for $F$ therefore gives
\[
 [F\det^n]\bigl(Z(YB)_{[1,a]}\bigr)
    =\left(\prod_{j<a}B_{jj}^{\nu_j+n}\right)
                     [F\det^n]\bigl(ZY_{[1,a]}\bigr).
\]
Taking coefficients in $Z$ gives precisely the target homogeneous
degrees and all target column-addition identities, separately at every
embedding.  Thus every $Q_\alpha(F)$ belongs to the original target
representation.  This is the actual $\operatorname{JetInVV}$ condition;
any target projection in the original definition acts as the identity.

The scalar in (3) is exactly
\[
 \eta_i(d,d')=
 \frac{\sigma\bigl(\mu_i^{-1}\!\det(N,\mathrm{fr})/
                          \mu_i^{-1}\!\det(N',\mathrm{fr}')\bigr)}
      {\chi_F(N,N')}.                                         \tag{4}
\]
The notation $\mu_i^{-1}\!\det$ denotes the original local
determinant-character factor.  To check (4), evaluate each such
factor on the specified frame: its character is the inverse of
the original finite Hecke character on the corresponding determinant
ideal.  This equality uses both its values on the chosen uniformizers
and its restrictions to local units.  Substituting it in the two
original $\rho$ factors cancels the intermediate determinant ideals
and leaves precisely (4).  The determinant multiplications in the
polynomial definition cancel the original frame shifts in the same
calculation.  Hence evaluation of (3) at the original class
representatives is literally the original $K_i$ kernel.  The original
Fourier average, positive index square root and period normalization
remain inside $J_\alpha$; this comparison adds no factor.

The original theta change-of-lattice identity, proved by reindexing
the normalized finite indicators, makes (3) equivariant under the
original marking groups.  Extend (3) complex linearly through
evaluation.  For arbitrary points represented as
$(b,\zeta,d)$ and $(b',\zeta',d')$, set
\[
 \widetilde k_i((b,\zeta,d),(b',\zeta',d'))
     =\sigma(\zeta')\,k_i(d,d')_{\mathbb C}\,\sigma(\zeta)^{-1}.
                                                               \tag{5}
\]
Equations (1) and equivariance of (3) prove that (5) is independent
of both representatives.  Formula (2) now gives source and target
equivariance.  At class representatives (5) is the literal $K_i$
kernel just identified.  This constructs the required point kernel;
there is no choice of a new projection or extension on a complement.

\emph{4. Local bounds and the Hecke identity.}
The established ordinary CM moment theorem for (3) has selected
support scalar $T^{-fN_{i-1}}$ and conjugate support scalar $1$.
In standard selected local frames, the actual source kernel lattice
has determinant $1$, and the target functional kernel has a basis
of determinant $T$.  Polynomial substitution therefore contributes
exactly $T^{-\sum_\tau s_{i-1,\tau}}$ in selected type.  In conjugate
type both dual lattices have determinant valuation $-1$ and the
difference of determinant shifts is exactly $h$, giving $T^{fh}$.
The direction scalars are the units furnished by the same ordinary
normalization.  The character ratio (4) is a unit as well: its ideal
arguments are prime to $T$, and a positive power of each algebraic
Hecke character has its original principal-ideal formula.  Raising
the ideal to its ray-class order reduces the assertion to a product
of powers of $T$-units.  Finite-order factors are units directly.

For completeness the passage from moments to the \emph{Weyl}
lattice is not left as a representation-theoretic hypothesis.
The established hook-product formula makes the normalized
highest-vector pairing a local unit under the original target hook
bound $\mu_1-\mu_{N_i}+N_i-1\le T$.  Indeed the residual partition
has last part zero, so its largest hook is at most
$\mu_1-\mu_{N_i}+N_i-2<T$; the endpoint introduces no hook equal
to $T$.  Its pairing is then integral
on the dual-Weyl lattice, proving equality with the Weyl lattice.
Integral substitution and simultaneous approximation of the finitely
many moments give exactly the selected and conjugate coefficient
kernel bounds in these identical lattices.  Unitary approximation
transports standard frames to all the actual local frames, using
the exact change-of-lattice identity already used for (5).
Step 2 transfers these bounds to the original point-value conditions.

Upper and lower neighbors in $P_i$ are the original neighbors: their
residue vectors satisfy the same equation $\ell(x)=1$, and their
other components are unchanged.  The finite matrix-incidence identity
for the same theta indicator is consequently
\[
 \sum_{p_1\ {
m above}\ p'}\widetilde k_i(p,p_1)
       =T^h\sum_{p_0\ {
m below}\ p}\widetilde k_i(p_0,p').
                                                               \tag{6}
\]
The equality follows on the rational polynomial evaluation basis
from the established indicator identity and hence on all values.
It includes its Hom-lattice covolume factor $T^h$.  Simultaneous
dilation of both residue functionals gives the same indicator and
proves the dilation law.  Distinct-prime neighbor operators commute,
since their local modifications have disjoint support.  Their ordered
product gives the commutation in \eqref{eq:actual-chain-step}.

Every class stabilizer is a finite subgroup of
$\operatorname{GL}_{2fN_i}(\mathbb Q)$.  If $T$ divided its order,
an element of order $T$ would have a primitive $T$-th root of unity
as an eigenvalue, forcing $T-1\le2fN_i$.  The retained strict bound
$T>2fN_m+1$ therefore makes each original mass denominator a local
unit.  Thus the finite orbit sum preserves both pointwise bounds.
Applying each selected neighbor removes its own allowance; removing
each remaining allowance costs exactly its residual first part
$r_{i,\tau}$.  This is the original integral weight decomposition
calculation, and yields the displayed $d_i$ and
\eqref{eq:actual-chain-step}.  In conjugate type retain the positive
factor $T^{fh}$ through the finite orbit sum and use the rational
image characterization of Step 2.  This proves the one-step part
of \eqref{eq:actual-chain-conjugate}.

\emph{5. The original initial source.}
At stage zero the reference character is trivial on the forgotten
extra markings.  Pullback along the actual forgetful class map
therefore leaves the underlying decorated lattice and its coefficient
value unchanged.  It preserves the original base integral lattice,
including every condition in $\mathscr S_{\rm ram}(w)$.  The
coefficient evaluation maps commute with this pullback, since the
frames and the polynomial are the same.  The original residue-dilation
norm operator $\mathrm{Nm}$ commutes as well: it is the same finite
sum over nonzero residue scalars, and every summand has the same
decorated lattice and value before and after the identification.
The upper neighbors also
commute with it: their unchanged underlying lattice and their equation
$\ell(x)=1$ give a bijection before and after forgetting the marking.
Thus the ordered product $U_0$, every iterate $U_0^m$, and the
prescribed scalar $T^a$ commute with the original initial
identification.  This proves precisely the original twisted source
tie, including its definition as a rational image, rather than
assuming it from isomorphism of ambient complex spaces.

Finally, multiply the integral maps $T^{d_i}U_iK_i$ in order.
Moving each $U_i$ to the right using
\eqref{eq:actual-chain-step} gives
\[
 (T^{d_m}U_mK_m)\cdots(T^{d_1}U_1K_1)
    =T^{\sum_i d_i+h|J|m(m+1)/2}C_mU_0^m.
\]
For an original rational-image member, divide by the already present
$T^a$ in that equality.  Step 2 gives the unique rational integral
target.  This proves \eqref{eq:actual-chain-selected}, also when
$\ell<0$.  Iterating the divided conjugate maps proves
\eqref{eq:actual-chain-conjugate}.  No inverse or nonvanishing of
any theta or neighbor map has been used.
\end{proof}

This discharges the selected-prime point-model, integral-coordinate,
kernel and twisted-source compatibility obligations for the actual
finite arithmetic theta chain.  It supplies no uniform bound for
the degree or discriminant of $K$.  The comparison of this chain
with the separately normalized global Poincar\'e period operators,
and their final retained source and target restrictions, is a further
identity; it is not inferred from agreement of representation types
or from the local estimates above.
